\subsection*{The mathematical facts}
\textbf{Three equivalent objects.} For each nonnegative integer n, these are in bijection: complete noncrossing pairings of 2n fixed points around a circle, with strings inside it (or in a row, with arcs above it); paths of n up-right steps U and n down-right steps D that start and end at height zero and never go below it; and finite ordered rooted trees with n edges. A tree is connected and has no cycles; the children at every vertex have a fixed left-to-right order. Circle labels stay fixed: rotations or reflections do not identify different partner lists.

\textbf{The maps are reversible.} Read pairing endpoints in order, writing U at a pair's first end and D at its other end. At each D, the inverse closes the most recent unclosed U; the noncrossing rule forces this choice. For a tree, visit whole child branches left-to-right, writing U when moving away from the root and D when returning. Its inverse adds a new rightmost child at U and moves to the parent at D. These rules prove the bijections; equal counts alone do not.

\textbf{The counts.} Let C(n) count any one of these sets and let C(0)=1. For n at least 1, C(n) is the sum of C(i)C(n-1-i) for i from 0 to n-1. The unique first-return split U P D Q proves this recurrence. Thus C(0),...,C(6) = 1, 1, 2, 5, 14, 42, 132. The full split argument is on page 9.

\textbf{Parity and neighbors.} An odd number of points cannot be completely paired. On a circle, a fixed chord can be completed without crossings exactly when it leaves an even number of other points on each side. Every nonempty complete noncrossing circle pairing contains a pair of circular neighbors. The proofs and exact instances follow on pages 3-4.

\subsection*{Where the worksheets investigate these facts}
\textbf{K-1 P1-P2:} Find the five six-dot pairings, then justify completeness. The two four-dot cases are printed examples.

\textbf{K-1 P3-P6:} Complete fixed chords; explore eight-dot pairings; test parity; decide whether neighbors can be avoided.

\textbf{Grades 2-3 P1-P4:} Enumerate six-dot pairings; encode paths; investigate the inverse; match all five paths and pairings.

\textbf{Grades 2-3 P5-P6:} Decide which codes are legal; enumerate all 14 eight-dot pairings.

\textbf{Grades 4-5 P1-P4:} Enumerate three-edge trees; encode and reconstruct trees; investigate uniqueness and characterize legal codes.

\textbf{Grades 4-5 P5-P6:} Count trees with four, five and six edges; propose and justify a general counting rule.

The examples introduce a representation or a walk, not a counting method. Let children conjecture before giving the facts above. Use drawings and explanations; keep hints and formal recurrence notation until needed.
